% rhode-island-thesis-doc.tex --- documentation for rhode-island-thesis.cls
%
% Copyright 2026 Kyle Monette
%
% This work may be distributed and/or modified under the
% conditions of the LaTeX Project Public License, either version 1.3c
% of this license or (at your option) any later version.
% The latest version of this license is in
%   https://www.latex-project.org/lppl.txt

\documentclass[11pt]{article}
\usepackage[T1]{fontenc}
\usepackage{lmodern}
\usepackage{amsmath}
\usepackage{booktabs}
\usepackage{tabularx}
\usepackage{geometry}
\geometry{margin=0.8in}
\usepackage{hyperref}
\hypersetup{colorlinks=true}
\setlength{\emergencystretch}{3em}

\newcommand\cls{\texttt{rhode-island-thesis}}
\newcommand\cs[1]{\texttt{\char`\\#1}}
\newcommand\meta[1]{\textit{$\langle$#1$\rangle$}}
\newcommand\marg[1]{\texttt{\{}\meta{#1}\texttt{\}}}
\newcommand\opt[1]{\texttt{#1}}
\newcommand\file[1]{\texttt{#1}}

\title{\texttt{rhode-island-thesis.cls}\\[0.5ex]
  \large Theses and dissertations at the University of Rhode Island}
\author{Kyle Monette (\href{https://kylemonette.github.io}{kylemonette.github.io}) \\
  \small Based on version 2.0 by Timothy Toolan \\
  \small Released under the LaTeX Project Public License (LPPL), version 1.3c}
\date{2026/10/28 --- Version 3.0}

\begin{document}
\maketitle

\section*{Overview}

\cls{} is a \LaTeX{} document class for writing a thesis or dissertation that
follows the formatting guidelines of the
\href{https://web.uri.edu/graduate-school/academics/thesis-dissertation/formatting/}{Graduate School}
at the University of Rhode Island.
It produces the title page, approval page, table of contents, list of figures,
and list of tables, and it sets the margins, spacing, and page numbering.
It also produces the List of References and the Bibliography from a Bib\TeX{}
database.

The package includes a sample thesis that compiles as-is.
Most students only need to edit it: fill in the details in \file{rhode-island-thesis.tex},
write chapters in \file{chapter1.tex}, \file{chapter2.tex}, and so on, and add
sources to \file{references.bib}.

This is a modified version of \file{urithesis.cls} version~2.0 by Timothy Toolan.
See Section~\ref{sec:license} for the license and the list of changes.

\tableofcontents

\section{Quick start}

\begin{enumerate}
\item Download the package and unzip it.
\item Compile \file{rhode-island-thesis.tex}:
  \begin{itemize}
  \item \textbf{Overleaf:} create a new project with \emph{New Project
    $\to$ Upload Project}, choose the zip file, and click \emph{Recompile}.
  \item \textbf{On your computer:} run \verb|latexmk -pdf rhode-island-thesis|, or run
    \verb|pdflatex rhode-island-thesis| followed by \verb|sh genbib.sh|
    (see Section~\ref{sec:compile}).
  \end{itemize}
\item Check that the result matches the sample \file{rhode-island-thesis.pdf}.
\item Edit \file{rhode-island-thesis.tex}: choose the class options
  (Section~\ref{sec:options}) and fill in the title, author, committee, and
  front matter (Section~\ref{sec:setup}).
\item Write your chapters and add your sources to \file{references.bib}.
\end{enumerate}

\section{Requirements}

\cls{} needs a \LaTeX{} kernel from October 2020 or later, pdf\LaTeX, and
Bib\TeX.
Any current \TeX{} Live, Mac\TeX, or MiK\TeX{} installation works, as does
Overleaf.
The class loads \texttt{etoolbox}, \texttt{ifthen}, and \texttt{url}.

\section{Files}

\noindent
\begin{tabularx}{\linewidth}{@{}>{\raggedright\arraybackslash}p{0.42\linewidth}X@{}}
\toprule
\file{rhode-island-thesis.cls} & The document class. \\
\file{uriref.bst} & Bib\TeX{} style for the numbered List of References, like [1]. \\
\file{uriapa.bst} & Bib\TeX{} style for author--year references, like (Toolan, 2006). \\
\file{uribib.bst} & Bib\TeX{} style for the alphabetical Bibliography. \\
\midrule
\file{rhode-island-thesis.tex} & The main file of the sample thesis. Compile this file only. \\
\file{abstract.tex}, \file{acknowledgements.tex}, \file{dedication.tex}, \file{preface.tex}
  & Front matter of the sample thesis. \\
\file{chapter1.tex}--\file{chapter3.tex}, \file{appendixA.tex}
  & Sample chapters and appendix, with examples of figures, tables,
    equations, and citations. \\
\file{references.bib} & Sample Bib\TeX{} database. \\
\file{rhode-island-thesis.pdf} & The compiled sample thesis. \\
\file{rhode-island-thesis-doc.tex}, \file{rhode-island-thesis-doc.pdf} & This document. \\
\bottomrule
\end{tabularx}

\medskip
\noindent Keep the class and \file{.bst} files in the same folder as
\file{rhode-island-thesis.tex}, unless your \TeX{} distribution already includes \cls.

\section{Class options}\label{sec:options}

Options go in the \cs{documentclass} line of \file{rhode-island-thesis.tex}, separated by
commas.
For example, a PhD dissertation with four committee members and page numbers
at the top right is written as
\begin{verbatim}
\documentclass[phd,4committee,topnum]{rhode-island-thesis}
\end{verbatim}

\subsection{General options}

These options all follow the Graduate School guidelines.
Defaults are marked with an asterisk.

\medskip
\noindent
\begin{tabularx}{\linewidth}{@{}lX@{}}
\toprule
\opt{masters}*, \opt{phd} & Master's thesis or doctoral dissertation. \\
\opt{12pt}*, \opt{11pt}, \opt{10pt} & Font size. \\
\opt{standard}*, \opt{manuscript} & Standard or manuscript format. In
  manuscript format, chapters are called ``Manuscript'' instead of
  ``Chapter''. \\
\opt{oneref}*, \opt{chapterref} & One List of References after the last
  chapter, or a List of References at the end of each chapter. \\
\opt{bibliography}, \opt{nobibliography} & Include or omit the alphabetical
  Bibliography at the end. The default is to include it in standard format,
  where the Graduate School requires it, and to omit it in manuscript format,
  where it is optional. \\
\opt{aparefs} & Cite by author and year, like (Toolan, 2006), instead of by
  number, like [1]. \\
\opt{3committee}*, \opt{4committee}, \opt{5committee} & Number of core
  committee members on the approval page, not counting the Dean. \\
\opt{sequential}*, \opt{nonsequential} & Number figures, tables, and
  equations throughout the thesis, or restart the numbering in each
  chapter. \\
\opt{bottomnum}*, \opt{topnum} & Page numbers at the bottom center or the top
  right. \\
\opt{hardcopy}*, \opt{electronic} & Approval page with signature lines, or
  with the typed names only. \\
\bottomrule
\end{tabularx}

\subsection{Drafting options}

These options make writing faster, but the result does not follow the
Graduate School guidelines.
Remove them before submitting your thesis.

\medskip
\noindent
\begin{tabularx}{\linewidth}{@{}lX@{}}
\toprule
\opt{draft} & Mark overfull lines with a black box, and show only the file
  names of figures. \\
\opt{draftbox} & Like \opt{draft}, but show the figures. \\
\opt{simpleref} & One List of References at the very end, and no
  Bibliography. Only one Bib\TeX{} run is needed. \\
\opt{noprelim} & Leave out the pages before Chapter~1. \\
\opt{nobib} & Leave out the List of References and the Bibliography.
  Citations appear as [?]. \\
\bottomrule
\end{tabularx}

\medskip
\noindent To compile only the chapter you are working on, combine
\opt{noprelim} and \opt{nobib} with \LaTeX's \cs{includeonly} command.

\section{Setting up \texttt{rhode-island-thesis.tex}}\label{sec:setup}

\subsection{Title page and approval page}

These commands go in the preamble of \file{rhode-island-thesis.tex}.

\medskip
\noindent
\begin{tabularx}{\linewidth}{@{}lX@{}}
\toprule
\cs{title}\marg{title} & Title of the thesis. Use the capitalization you want
  on the title page. \\
\cs{author}\marg{name} & Your name. Write \texttt{John L.\string~Doe} so that the
  period after an initial is not treated as the end of a sentence. \\
\cs{dept}\marg{department} & Department on the title page. The default is
  Electrical Engineering, so most students need to set this. \\
\cs{copyrightyear}\marg{year} & Year on the title and approval pages. The
  default is the current year; the Graduate School requires the year of your
  graduation. \\
\cs{signature}\marg{name} & One committee member. Use one \cs{signature} per
  member; the first is the major professor. The number must match the
  \opt{3committee}, \opt{4committee}, or \opt{5committee} option. \\
\cs{deansignature}\marg{name} & The Dean of the Graduate School. \\
\bottomrule
\end{tabularx}

\subsection{Front matter}

The title page, approval page, table of contents, list of figures, and list of
tables are produced automatically.
The lists of figures and tables appear only if the thesis has figures or
tables.
The other front matter is written in separate files, which are included with
the commands below.
Each argument is a file name without the \file{.tex} extension, and each file
contains only the text, since the heading is added automatically.

\medskip
\noindent
\begin{tabularx}{\linewidth}{@{}llX@{}}
\toprule
Section & Command & \\
\midrule
Abstract & \cs{abstract}\marg{file} & Required. \\
Acknowledgments & \cs{acknowledgements}\marg{file} & Optional. \\
Dedication & \cs{dedication}\marg{file} & Optional. \\
Preface & \cs{preface}\marg{file} & Optional. \\
\bottomrule
\end{tabularx}

\subsection{Chapters and appendices}

After \verb|\begin{document}|, include each chapter with
\cs{newchapter}\marg{file} and each appendix with \cs{newappendix}\marg{file}:
\begin{verbatim}
\newchapter{chapter1}
\newchapter{chapter2}
\newappendix{appendixA}
\end{verbatim}
Each chapter or appendix file starts with \cs{chapter}\marg{title}.
Chapters are numbered 1, 2, \ldots{} and appendices A, B, \ldots.
Do not put \cs{chapter} directly in \file{rhode-island-thesis.tex}.

\section{References and the Bibliography}

The Graduate School recognizes two kinds of reference lists.
\begin{description}
\item[List of References.] A numbered list of the sources cited, in the order
  they are cited. It is optional, and appears either at the end of each
  chapter (\opt{chapterref}) or once after the last chapter (\opt{oneref}).
\item[Bibliography.] An alphabetical list of all sources used or consulted.
  It is the last section of the thesis. It is required in standard format and
  optional in manuscript format.
\end{description}

\noindent \cls{} produces both from the same Bib\TeX{} database.
Name the database with \cs{reffile}\marg{file}, without the \file{.bib}
extension; a comma-separated list of files is also accepted.
Cite sources with \cs{cite} as usual.
Do not use \cs{bibliography} or \cs{bibliographystyle}; the class handles
both, and these commands do nothing.

To list a source in the Bibliography without citing it, use
\cs{nocite}\marg{key}, or \verb|\nocite{*}| to list every entry in the
database.
\cs{nocite} adds sources to the Bibliography only, never to the List of
References.

The styles are based on \texttt{IEEEtran.bst}, so its entry types are
available, including \texttt{@electronic} for web pages and \texttt{@patent},
\texttt{@periodical}, and \texttt{@standard}.
They are described in ``How to Use the IEEEtran Bib\TeX{} Style'' by Michael
Shell (run \verb|texdoc IEEEtran_bst_HOWTO|).
With \opt{aparefs}, an entry without an author is cited by its \texttt{key}
field.

\section{Compiling}\label{sec:compile}

A thesis has several reference lists, and each one needs its own Bib\TeX{}
run.
How you compile depends on your setup.
\begin{description}
\item[Overleaf.] Click \emph{Recompile}. Overleaf runs Bib\TeX{} on every list
  automatically.
\item[latexmk.] Run \verb|latexmk -pdf rhode-island-thesis|. It runs pdf\LaTeX{} and
  Bib\TeX{} as many times as needed. Many editors can use latexmk.
\item[By hand.] Run \verb|pdflatex rhode-island-thesis| once, then \verb|sh genbib.sh|.
  On every run, \cls{} writes the script \file{genbib.sh}, which runs Bib\TeX{}
  on each list and then pdf\LaTeX{} twice more. It also writes
  \file{genbib.txt}, which lists only the Bib\TeX{} commands, for systems
  without \texttt{sh}. Run these commands, then run \verb|pdflatex rhode-island-thesis|
  twice.
\end{description}

\section{Tips}

\begin{itemize}
\item The text is double spaced. Use \cs{singlespace} and \cs{doublespace} to
  switch spacing where the Graduate School allows single spacing, such as long
  quotations.
\item Load \texttt{hyperref} to make citations, cross-references, and the
  table of contents clickable in the PDF.
\item Do not load the \texttt{chapterbib} package. \cls{} contains its own
  version of that code, and the two conflict.
\item Figures and tables in appendices are not listed in the List of Figures
  or the List of Tables.
\end{itemize}

\section{License and changes}\label{sec:license}

\subsection{License}

This work may be distributed and/or modified under the conditions of the
\LaTeX{} Project Public License, either version 1.3c of this license or (at
your option) any later version.
The latest version of this license is at
\url{https://www.latex-project.org/lppl.txt}.

\cls{} is a modified version of \file{urithesis.cls} version~2.0 (2016),
copyright 2004--2016 University of Rhode Island, written by Timothy Toolan.
Timothy Toolan does not support this modified version.
Report problems to Kyle Monette at
\url{https://github.com/kylemonette/thesis-template/issues}.
The unmodified version~2.0 is available from
\url{https://www.ele.uri.edu/info/thesis/dist/thesis.zip}
and is archived at
\url{https://web.archive.org/web/20260216180715/https://www.ele.uri.edu/info/thesis/dist/thesis.zip}.

The Bib\TeX{} styles are distributed under their own licenses, which are
stated in each file.
\file{uriref.bst} and \file{uribib.bst} are based on \texttt{IEEEtran.bst} and
\texttt{IEEEtranS.bst} by Michael Shell (Perl Artistic License), and
\file{uriapa.bst} is based on \texttt{apalike.bst} by Oren Patashnik.

\subsection{Changes from version 2.0}

\paragraph{\file{rhode-island-thesis.cls}}
\begin{itemize}
\item The class was renamed from \file{urithesis.cls} to
  \file{rhode-island-thesis.cls}. Load it with
  \verb|\documentclass{rhode-island-thesis}|.
\item Requires a \LaTeX{} kernel from October 2020 or later.
  Options are stored with \texttt{etoolbox} booleans instead of \cs{newif}.
\item The default List of References is now one list after the last chapter
  (\opt{oneref}). Version~2.0 put a list at the end of each chapter by
  default; this is now the \opt{chapterref} option.
\item New options \opt{bibliography} and \opt{nobibliography} control the
  Bibliography. It is included by default in standard format and omitted by
  default in manuscript format. Version~2.0 always included it, except with
  \opt{simpleref} or \opt{nobib}.
\item The \opt{standard} option now works. In version~2.0 it referred to an
  undefined command.
\item The class writes \file{genbib.sh}, a script with the full compile
  sequence, in addition to \file{genbib.txt}. Its output files are opened with
  \cs{newwrite} instead of fixed stream numbers.
\item The redefinition of \cs{@include} uses the \LaTeX{} 2020 include hooks,
  so packages such as \texttt{hyperref} work with \cs{newchapter}.
\item The end-of-document code uses \cs{AtEndDocument} instead of redefining
  \cs{enddocument}.
\item The per-chapter counter reset list (\cs{cl@chapter}) is handled so that
  packages such as \texttt{listings} that extend it do not break the class.
\item The comments and option descriptions were rewritten, and the code was
  reorganized. For a line-by-line record, compare this file with version~2.0.
\end{itemize}

\paragraph{\file{uriref.bst}, \file{uribib.bst}, \file{uriapa.bst}}
Unchanged, except for a notice at the top of each file and the class name in
the comments.

\paragraph{Sample files}
\file{thesis.tex} was renamed \file{rhode-island-thesis.tex}.
It, \file{abstract.tex}, \file{chapter1.tex}, and
\file{references.bib} were rewritten.
\file{chapter2.tex}, \file{chapter3.tex}, \file{appendixA.tex},
\file{acknowledgements.tex}, \file{dedication.tex}, and \file{preface.tex} are
new.
\file{thesisbib.tex} was removed; it is no longer needed.
This document is new.

\end{document}
